Además de los módulos de investigación específicos, se propone implementar un servicio de asistencia técnica orientado a acompañar al equipo de YPF en el análisis de preguntas que surjan en el curso ordinario de la operación y la toma de decisiones. La idea no es reemplazar las capacidades analíticas internas, sino complementarlas: aportar una perspectiva independiente y capacidad de ejecución en situaciones donde el equipo interno enfrenta restricciones de tiempo, expertise específico o volumen de trabajo.

\subsection{Alcance}

Este servicio esta pensado para cubrir tres tipos de intervenciones:

\begin{itemize}

    \item \textbf{Análisis estadísticos puntuales.} Procesamiento y análisis de datos sobre variables clave del negocio —precios, márgenes, volúmenes de despacho, participación de mercado— en respuesta a preguntas concretas que requieran una respuesta rápida y fundada.

    \item \textbf{Evaluación de políticas y campañas.} Estimación del impacto de intervenciones comerciales u operativas de alcance acotado —una promoción,  un cambio de precios, una campaña regional— que no justifiquen por sí solas  un módulo de investigación completo pero sí requieren una evaluación más  rigurosa que una comparación descriptiva simple.

    \item \textbf{Soporte metodológico.} Asistencia al equipo interno en el diseño de mediciones, la revisión de metodologías existentes, la interpretación de resultados o la identificación de limitaciones en análisis en curso.

\end{itemize}
